\begin{explanationbox}
	\textbf{\textcolor{CExplanation}{一句话直觉}}

	圆心角张开的度数就是所对弧在圆上划过的度数，两者数值完全相同。

	\medskip
	\textbf{\textcolor{CExplanation}{核心拆解}}
	\begin{description}[style=nextline, leftmargin=2.8em, labelsep=0.5em, itemsep=0.45em]
		\item[\textbf{要点一（定义对应）：}] 圆心角的顶点在圆心，两边是半径；所对弧是两边之间的圆周部分。这种对应是唯一的。
		\item[\textbf{要点二（度数等同）：}] 圆心角与所对弧的度数完全一致，因为两者刻画的是同一段旋转量。
	\end{description}

	\medskip
	\textbf{\textcolor{CExplanation}{几何本质（旋转量对应）}}

	圆心角是半径从$OA$旋转到$OB$的转动角度，而弧$\widehat{AB}$是此旋转过程中点$A$沿圆周运动的轨迹。转过的角度就是弧$\widehat{AB}$对应的中心角度数，因此两者相等。

	\medskip
	\textbf{\textcolor{CExplanation}{代数意义（数值转换）}}

	在数值上，圆心角的度数等于所对弧的度数。因此，在计算或推理中，可以直接用弧的度数替换圆心角的度数，反之亦然。

	\medskip
	\textbf{\textcolor{CExplanation}{考点价值（命题方向）}}
	\begin{itemize}[leftmargin=*, itemsep=0.5em, topsep=0.4em]
		\item \textbf{考法一（直接计算）：} 给出圆心角求所对弧的度数，或给出弧的度数求圆心角。
		\item \textbf{考法二（等量替换）：} 利用圆心角相等证明所对弧相等（或反之），解决圆中角与弧的转化问题。
	\end{itemize}

	\medskip
	\textbf{\textcolor{CExplanation}{顿悟点}}

	圆心角与所对弧本质是“同一个量”的两种表示：旋转角和轨迹弧。抓住了这一点，圆中的角度与弧的换算就变成了简单的一一对应。

	\medskip
	\textbf{\textcolor{CExplanation}{使用场景}}
	\begin{itemize}[leftmargin=*, itemsep=0.5em, topsep=0.4em]
		\item \textbf{场景一（单一计算）：} 在圆中直接给出圆心角或弧的度数，求另一个。
		\item \textbf{场景二（等量推理）：} 结合圆心角相等、弧相等的关系，进行几何推理。
	\end{itemize}
\end{explanationbox}
